\documentclass[11pt]{article}

\usepackage[a4paper,margin=2.5cm]{geometry}
\usepackage[utf8]{inputenc}
\usepackage[T1]{fontenc}
\usepackage{lmodern}
\usepackage{microtype}
\usepackage{hyperref}
\usepackage{xcolor}
\usepackage{listings}
\usepackage{booktabs}
\usepackage{enumitem}
\usepackage{parskip}
\usepackage{ragged2e}
\emergencystretch=2em
\sloppy

\hypersetup{
  colorlinks=true,
  linkcolor=black,
  urlcolor=blue,
  citecolor=black
}

\definecolor{codebg}{RGB}{246,246,246}
\lstset{
  basicstyle=\ttfamily\small,
  backgroundcolor=\color{codebg},
  frame=single,
  framerule=0.3pt,
  rulecolor=\color{gray!40},
  breaklines=true,
  columns=fullflexible,
  showstringspaces=false,
  tabsize=2,
  xleftmargin=8pt,
  xrightmargin=8pt
}

\title{From UIMA CAS to a Video Player:\\ How the Bundestag Pipeline Fits Together}
\author{Video-Emotion Recognition project notes}
\date{\today}

\begin{document}
\maketitle

\section{What this document is for}

The Java DUUI pipeline that actually runs face detection, speaker diarization,
transcription and the three emotion models is not part of this repository ---
it lives elsewhere and hands us a single artifact: a UIMA CAS serialized as an
\texttt{.xmi} file, together with the source video. Everything described here
starts from that file. There are three moving parts: a Python parser that
turns the CAS into rows in Postgres, the Postgres schema itself, and a small
FastAPI + vanilla-JS viewer that plays the video back with everything overlaid
in sync. This note walks through all three in the order data actually flows
through them.

\section{Why the CAS is annoying to parse}

A UIMA CAS cannot be read without the typesystem that describes it --- an XML
document listing every annotation type and its fields. The three typesystem
files we were given
(\texttt{IdentityEmotionTypeSystem.xml},
\texttt{MultimodalIdentityTypeSystem.xml},
\texttt{EmotionTypeSystem.xml}) turned out to be mostly copies of each other
rather than complementary pieces, and between them they still reference a
handful of types --- \texttt{MultimediaElement}, \texttt{Embedding}, the bare
\texttt{MetaData}, \texttt{AudioToken} --- that none of the three actually
define. Those come from external shared typesystem imports that never shipped
with the export.

\texttt{duui\_parser/typesystem.py} deals with this once, up front:
\texttt{load\_merged\_typesystem()} concatenates the three files'
type descriptions, drops exact duplicates by name, and injects hand-written
fallback definitions for the types that are referenced but missing, based on
what their fields actually look like on real annotation instances. Anything
still undefined after that gets a generic featureless stub so loading
succeeds instead of raising.

Loading the CAS itself
(\texttt{pipeline.py}) uses two non-default flags:

\begin{lstlisting}[language=Python]
cas = load_cas_from_xmi(f, typesystem=merged_typesystem,
                         lenient=True, trusted=True)
\end{lstlisting}

\texttt{lenient=True} means an annotation whose type is not in the merged
typesystem at all --- the DKPro linguistic layer (Token, Lemma, POS,
Dependency, NamedEntity), which this pipeline never populates in a form we
use --- is silently skipped rather than aborting the whole load.
\texttt{trusted=True} switches on lxml's huge-tree parsing; a full-session
export is a single-line XML document north of 40\,MB and lxml's default
buffer limit rejects it otherwise.

\section{The parser}

\texttt{duui\_parser/parsers/} has one module per table family, and they run
in a fixed order (\texttt{parsers/\_\_init\_\_.py}):

\begin{lstlisting}[language=Python]
PARSE_STEPS = [
    video, model, person, segment, token,
    embedding, presence, detection, emotion, text_emotion,
]
\end{lstlisting}

The order is not arbitrary. \texttt{video} runs first because everything
else needs the video's id. \texttt{person} runs early because the schema's
foreign keys mean a \texttt{persons} row has to exist before anything that
references \texttt{person\_id} can be inserted, and because it also builds
two lookup dictionaries --- \texttt{face\_id\_to\_person\_id} and
\texttt{voice\_id\_to\_person\_id} --- that every later step needs.

Every module follows the same shape: walk every UIMA view with
\texttt{select\_across\_views} (video-derived and transcript-derived
annotations sit on separate sofas in this pipeline, so a plain
\texttt{cas.select()} on one view would miss half the data), then run an
\texttt{INSERT ... ON CONFLICT DO NOTHING/UPDATE} per instance. The XMI id
of each annotation is reused directly as its Postgres primary key, which is
what makes re-running the same file idempotent --- rerun it twice and you get
the same rows, not duplicates. The whole import happens inside one
transaction; anything raising partway through rolls the entire thing back,
so a half-imported video should never exist.

\subsection{Identity resolution}

This is the part that took the most debugging. \texttt{FaceIdentity} and
\texttt{VoiceIdentity} annotations each carry a bare string id
(\texttt{faceId="face\_1"}, \texttt{voiceId="SPEAKER\_00"}). A separate
\texttt{identity:Person} annotation is the only thing that actually links
those strings to a person, via a pipe-delimited label:

\begin{lstlisting}
<identity:Person xmi:id="32539" personId="person_1"
                  label="voice:SPEAKER_00|face:face_1|score:0.83"/>
\end{lstlisting}

\texttt{identity\_resolution.py} parses that label and, for every
downstream annotation that only carries a bare \texttt{faceId}/\texttt{voiceId}
(detections, embeddings, presences, emotions), looks it up against the maps
built from \texttt{Person} rows. If a given CAS export has no
\texttt{identity:Person} annotations at all --- which happened with one of
the two full-session exports we were given, and not with the other --- every
one of those lookups returns \texttt{None} for the whole file, regardless of
how many \texttt{FaceIdentity}/\texttt{VoiceIdentity} records exist. That is
a property of the source export, not a parser bug; the fix was getting a
different \texttt{.xmi} that actually had the identity-linking stage run.

The important guarantee here, worth stating explicitly because it is easy to
assume otherwise: \emph{a failed identity lookup never drops the row}.
Every detection, presence interval and emotion reading gets inserted with
whatever \texttt{person\_id} it resolved to, including \texttt{NULL}. On the
video currently loaded, that is the difference between 1970 face detections
total and 1474 of them actually tagged with a person --- the other 496 are
real, present rows for faces the pipeline saw but never linked to an
identity.

\section{The database}

\texttt{schema/schema.sql} defines fifteen tables, which fall into three
groups.

\textbf{Metadata:} \texttt{videos} (one row per source file) and
\texttt{models} (provenance --- which face/voice/emotion model produced a
given row, since several are in play).

\textbf{Identity:} \texttt{global\_persons} (cross-video identity, almost
never populated in practice --- cross-video clustering is the exception) and
\texttt{persons} (per-video identity, the one that actually matters day to
day), backed by \texttt{face\_embeddings}/\texttt{voice\_embeddings}
(native \texttt{pgvector} columns, 512-dim and 192-dim), with
\texttt{presences} recording on-screen/speaking intervals and
\texttt{face\_detections}/\texttt{person\_detections} carrying the actual
per-frame bounding boxes, normalized to $[0,1]$ so they are resolution
independent.

\textbf{Content and emotion:} \texttt{segments} covers both camera shots and
spoken sentences and, for sentences, carries both a time window and a
character-offset window into the transcript --- the bridge that lets you
map a moment in the video to a span of text. \texttt{linguistic\_tokens} is
one row per spoken word; sentences do not store their own text, so the
frontend reconstructs a subtitle by joining every token whose
\texttt{start\_time} falls inside the sentence's window, deliberately not by
foreign key, because the token-to-segment FK is not reliably populated by
the real annotator. \texttt{base\_emotions} is the central table: one row
per emotion reading, tagged by \texttt{modality} (video / audio / text) and
\texttt{granularity}, with a genuinely different set of populated columns
per modality --- video rows are per-frame with a real bounding box, audio
rows carry a dominance score that text rows never have, and text rows in
practice come from two different upstream annotators of differing
completeness. \texttt{emotion\_scores} holds the full label distribution
behind each reading's \texttt{dominant\_label}, and \texttt{fused\_emotions}
holds an optional cross-modal aggregate that is empty for most runs,
including the one currently loaded.

A theme worth naming: nearly every foreign key in this schema is nullable
on purpose, because real annotator output is incomplete in ways the schema
has to tolerate rather than hide. Most \texttt{presences} rows have no
resolved \texttt{person\_id}. \texttt{fused\_emotions} is routinely empty.
Text-modality rows sometimes have no \texttt{start\_time} at all. None of
this is a defect to fix; it is what the data actually looks like, and every
consumer downstream --- the frontend, the query agent --- is written to
expect it.

\section{The frontend}

\texttt{webapp/server.py} is deliberately thin. It exposes three things: a
list of videos, one JSON payload per video containing every sentence, shot,
emotion reading and detection for it, and the raw video bytes themselves
served with range-request support so seeking works. The client fetches that
one payload exactly once per video and never calls the API again during
playback --- everything after that is a client-side lookup against
in-memory arrays.

\texttt{webapp/static/app.js} hooks the video element's \texttt{timeupdate}
event when paused and \texttt{requestAnimationFrame} while playing, and on
every tick asks two questions of the data it already has:

\begin{lstlisting}
// which rows genuinely span this instant
function coveredBy(rows, t) {
  return rows.filter(r => r.start_time <= t && t <= r.end_time);
}

// which frame-sampled row is closest, within tolerance
function nearestByTime(rows, t, tolerance) {
  // detections are sampled per frame, not continuous,
  // so "covers this instant" doesn't apply to them
}
\end{lstlisting}

Sentences and time-anchored emotions use the first kind of lookup; face and
person detections use the second, since they are sampled at discrete frames
rather than spanning a continuous interval. Boxes are drawn on a canvas
overlaid on the video element, scaled from the stored $[0,1]$ coordinates to
whatever size the video is actually rendered at, colored by
\texttt{person\_id} through a small fixed palette (grey for anything
unresolved), and tagged with a video-emotion label by joining on the shared
\texttt{frame\_index} between a detection and the emotion reading computed
from it.

A natural-language ``Ask'' panel sits above the player and is worth
mentioning separately because it is the one piece that talks to an external
model rather than just Postgres. It sends the question to Claude along with
a hand-written schema description --- the semantics above, not just column
names, since things like ``compare emotions across modalities by valence
rather than by label string'' are not recoverable from
\texttt{information\_schema} alone --- Claude proposes and previews a SQL
query through a tool-use loop, and only the final validated query is
actually executed, inside a read-only transaction with a statement timeout
and a row cap, never trusting rows the model claims to have already seen.

\section{Putting it together}

Concretely, going from a fresh \texttt{.xmi} to a working demo is:
\texttt{python main.py cas/your\_file.xmi} to parse and commit, then
\texttt{uvicorn webapp.server:app --reload} to serve it. Nothing is ever
baked into the video file itself --- every overlay is computed at playback
time from what is in Postgres, so improving a detection model, re-running
identity resolution, or fixing a mislabeled emotion is always just a matter
of re-parsing and reloading the page, never re-encoding video.

\end{document}
